\section{Consequences of the proof}\label{sec:consequences}
The results of this section are not needed for Theorem~\ref{thm:main}. They make the
construction quantitative and more flexible.

\subsection{The feedback deficit in closed form}
\begin{proposition}\label{prop:deficit}
For every \(n\ge2\), the stability margin \(1-G_n\) of the limiting prefix in
Lemma~\ref{lem:gain} is
\begin{equation}\label{eq:deficit}
 1-G_n=\frac{(3n-2)(50n-45)\,(2/9)^{n-1}}{9\,\bigl[14-8(3n-2)(2/9)^{n-1}\bigr]},
 \qquad 1-G_n\sim\frac{25}{21}\,n^2\Bigl(\frac29\Bigr)^{n-1}\quad(n\to\infty).
\end{equation}
\end{proposition}
\begin{proof}
Insert \(v=(3n-2)(2/9)^{n-1}\) from \eqref{eq:vmu} into \eqref{eq:gainfinal}; the asymptotic
statement follows from \(v\to0\) and \(3\cdot50/(9\cdot14)=25/21\).
\end{proof}

\begin{table}[ht]
\centering
\caption{The coefficient fraction \(v_n\) and the feedback deficit \(1-G_n\) of the limiting
prefix (exact values and decimal approximations).}\label{tab:deficit}
\begin{tabular}{cccccccc}
\toprule
\(n\) & 2 & 3 & 4 & 5 & 6 & 7 & 8\\
\midrule
\(v_n\) & \(\frac89\) & \(\frac{28}{81}\) & \(\frac{80}{729}\) & \(\frac{208}{6561}\) &
 \(\frac{512}{59049}\) & \(\frac{1216}{531441}\) & \(\frac{2816}{4782969}\)\\[4pt]
\(1-G_n\) & \(\frac{220}{279}\) & \(\frac{14}{39}\) & \(\frac{6200}{43047}\) &
 \(\frac{4264}{81171}\) & \(\frac{4352}{246777}\) & \(\frac{185440}{33437007}\) &
 \(\frac{499840}{301225671}\)\\[4pt]
decimal & \(0.789\) & \(0.359\) & \(0.144\) & \(0.0525\) & \(0.0176\) & \(0.00555\) & \(0.00166\)\\
\bottomrule
\end{tabular}
\end{table}

\begin{figure}[ht]
\centering
\includegraphics[width=\linewidth]{figures/deficit.pdf}
\caption{(a) The feedback deficit \eqref{eq:deficit}, its asymptote, and the coefficient
fraction \(v_n\). (b) The entrywise distance between \(r\bA_r\), computed exactly from the
coalesced construction, and its limit \(\bB_*\); the slope is \(-1\), in agreement with the
\(O(1/r)\) errors in Section~\ref{sec:prefix}.}\label{fig:deficit}
\end{figure}

The deficit is positive for every \(n\) but decays geometrically (Table~\ref{tab:deficit},
Figure~\ref{fig:deficit}(a)). This explains why the construction works for all \(n\) and is
nevertheless poorly conditioned for large \(n\). The decay is a property of this construction,
with all roots coalesced at \(x=3\) and \(A=(1+u)D\); it is not a lower bound on the time
scales or the resources that \(n\) stable states require.

\subsection{Diagonal stability and independent catalytic currents}
Lemma~\ref{lem:gain} used one vector of catalytic currents, \((1,\dots,1,t)\). A diagonal
Lyapunov matrix removes this restriction. Let \(\bz=\bH^{-1}\one\) and \(\bw=\bH^{-1}\bphi\);
both are entrywise positive.

\begin{proposition}[Diagonal Lyapunov matrix]\label{prop:diagonal}
Let \(\bP=\diag(w_i/z_i)\) and \(\bQ_*=\bB_*^{T}\bP+\bP\bB_*\). Then \(\bQ_*\) is negative
definite. Consequently, for all sufficiently large \(r\), \(\bA_r^{T}\bP+\bP\bA_r\) is negative
definite, and for every such \(r\) and every positive diagonal matrix \(\boldsymbol\Omega\) the
matrix \(\bA_r\boldsymbol\Omega\) is Hurwitz.
\end{proposition}
\begin{proof}
Since \(\bH\) is symmetric, \(\one^{T}\bw=\bz^{T}\bphi=G\). Hence \(\bB_*\bz=-(1-G)\one\) and
\(\bB_*^{T}\bw=-\bphi+G\bphi=-(1-G)\bphi\). Because \(\bP\bz=\bw\),
\[
 \bQ_*\bz=\bB_*^{T}\bw+\bP\bB_*\bz=-(1-G)\bigl(\bphi+\bP\one\bigr)<0 .
\]
The matrix \(\bQ_*\) is symmetric with nonnegative off-diagonal entries \(Q_{ik}\), because
\(\bB_*\) has nonnegative off-diagonal entries and \(\bP\) is a positive diagonal matrix. For
such a matrix and any positive vector \(\bz\), expanding the squares gives the identity
\begin{equation}\label{eq:energy}
 \boldsymbol a^{T}\bQ_*\boldsymbol a=\sum_i\frac{(\bQ_*\bz)_i}{z_i}\,a_i^2
 -\sum_{i<k}Q_{ik}\,z_iz_k\Bigl(\frac{a_i}{z_i}-\frac{a_k}{z_k}\Bigr)^2 ,
\end{equation}
which is negative for \(\boldsymbol a\ne0\). Negative definiteness is an open condition, and
\(r\bA_r\to\bB_*\). Finally
\((\bA_r\boldsymbol\Omega)^{T}(\boldsymbol\Omega\bP)+(\boldsymbol\Omega\bP)(\bA_r\boldsymbol\Omega)
=\boldsymbol\Omega(\bA_r^{T}\bP+\bP\bA_r)\boldsymbol\Omega\) is negative definite, and
\(\boldsymbol\Omega\bP\) is positive definite, so \(\bA_r\boldsymbol\Omega\) is Hurwitz by
Lyapunov's theorem.
\end{proof}

\begin{corollary}[Independent prefix currents]\label{cor:currents}
Let \(n\ge2\), let \(r\) be as in Proposition~\ref{prop:diagonal}, and let
\(\mathcal K\subset\R^{n-1}_{>0}\) be compact. There are \(t>0\), \(\varepsilon>0\) and
\(\delta_0>0\) such that the conclusions of Theorems~\ref{thm:main}\textup{(b)}
and~\ref{thm:split} hold for the construction of Section~\ref{sec:lift} with the catalytic
currents \((\omega_1,\dots,\omega_{n-1},t)\) at the coalesced equilibrium, for every
\((\omega_1,\dots,\omega_{n-1})\in\mathcal K\) and every \(0<\delta\le\delta_0\).
\end{corollary}
\begin{proof}
By Proposition~\ref{prop:diagonal} the prefix \(\bA_r\diag(\omega_1,\dots,\omega_{n-1})\) is
Hurwitz for every point of \(\mathcal K\), with a spectral bound that is uniform by
compactness. The arguments of Lemmas~\ref{lem:slowlink} and~\ref{lem:fullspectral} and of the
proof of Theorem~\ref{thm:split} use only the continuity of eigenvalues and the exact zero
provided by \eqref{eq:detzero} and Lemma~\ref{lem:Ksingular}. All matrices involved depend
continuously on the point of \(\mathcal K\), so \(t\), \(\varepsilon\) and \(\delta_0\) can be
chosen uniformly on \(\mathcal K\).
\end{proof}

The currents \(\omega_1,\dots,\omega_{n-1}\) are source-preserving coordinates; they are not
independent perturbations of the original rate constants, which are covered by
Theorem~\ref{thm:main}(c). No uniformity over unbounded sets of currents is claimed.

\subsection{Readout and the separation--recovery trade-off}
A natural observable is the total amount of fully phosphorylated substrate, \(S_n+Y_n\). At the
equilibrium of the construction of Section~\ref{sec:construction} with auxiliary coordinate
\(x\), formulas \eqref{eq:states} and \eqref{eq:sf} with \(t_n=D_{n-1}\) give
\begin{equation}\label{eq:readout}
 S_n+Y_n=\mathcal R(x):=\frac{D_{n-1}}{2}\cdot\frac{u^{n-1}}{D(u)}\cdot\frac{(x-1)(x+5)}{x+1},
 \qquad u=\frac{x^2-1}{8}.
\end{equation}
The function \(u^{n-1}/D(u)=1/\sum_iD_iu^{i-(n-1)}\) is nondecreasing, and
\((x-1)(x+5)/(x+1)=x+3-8/(x+1)\) has derivative \(1+8/(x+1)^2>0\). Hence \(\mathcal R\) is
strictly increasing: the readout orders the \(2n-1\) equilibria in the same way as the ratio
\(E/F\), and sinks and saddles alternate along it. Also \(E=(x-1)/2\) and \(F=4/(x+1)\) at these
equilibria, so the total phosphorylation flux \(\sum_i\gamma_iY_i\) is positive at each of
them; they are nonequilibrium steady states sustained by the irreversible catalytic steps.

\begin{corollary}\label{cor:tradeoff}
In the split family of Theorem~\ref{thm:split},
\[
 \mathcal R(x_j)-\mathcal R(x_k)=\mathcal R_\circ'(3)\,(\xi_j-\xi_k)\,\delta+O(\delta^2),\qquad
 \frac{1}{|\lambda_j(\delta)|}=\frac{1+O(\delta)}{a\,\prod_{k\ne j}|\xi_j-\xi_k|}\;\delta^{-(2n-2)},
\]
where \(\mathcal R_\circ\) is the function \eqref{eq:readout} of the coalesced system and
\(\mathcal R_\circ'(3)>0\).
\end{corollary}
\begin{proof}
At each \(\delta\) all equilibria share the polynomial \(D=D^{(\delta)}\), so the dependence of
\(D\) on \(\delta\) cancels from the difference to first order. The second formula is
\eqref{eq:splitlaw}.
\end{proof}

Thus, within this near-coalescence family, separating the states by an amount \(\delta\) costs
a linear relaxation time of order \(\delta^{-(2n-2)}\). This is a statement about the family,
not a speed limit for the network, and it suggests looking for well-separated roots when
designing examples; Section~\ref{sec:benchmark} does so.

\subsection{Labels, and a terminating design procedure}
\begin{corollary}\label{cor:labels}
A single sequential distributive cycle admits \(q\) simultaneously stable positive equilibria,
for one rate vector and one triple of totals, if and only if \(n\ge q\). In particular
\(2^{b}\) distinct stable equilibrium labels require \(n\ge2^{b}\) sites.
\end{corollary}
\begin{proof}
If \(n\ge q\), Theorem~\ref{thm:main}(b) gives \(n\ge q\) sinks; the converse is
Theorem~\ref{thm:main}(a).
\end{proof}

Coupled modules, transient codes and oscillatory phases are outside this statement, and it is
not a statement about an information rate in the presence of noise.

\begin{remark}[Exact design]\label{rem:design}
For fixed \(n\) the choices \eqref{eq:order} can be made by a terminating procedure in rational
arithmetic: double \(r\) until \(\bA_r\) passes an exact Routh test (or until
\(-(\bA_r^{T}\bP+\bP\bA_r)\) passes an exact \(LDL^{T}\) test); halve \(t\) until the slow
matrix has a simple zero and a Hurwitz deflation; halve \(\varepsilon\) until the coalesced
restricted Jacobian has the same property; halve \(\delta\) until all \(2n-1\) exact Routh
tests give the alternating pattern. Each loop terminates by the corresponding lemma. No bound
on the number of steps is claimed. Proposition~\ref{prop:witness} reports the outcome for
\(n\le7\).
\end{remark}

\section{Exact examples}\label{sec:examples}

\subsection{Rational witnesses for \texorpdfstring{\(n\le7\)}{n at most 7}}
The proof of Theorem~\ref{thm:main}(b) is asymptotic in \(r\), \(t\), \(\varepsilon\) and
\(\delta\). For small \(n\) the procedure of Remark~\ref{rem:design} stops at once: no slow
final link and no fast binding are needed, and \(r\) and \(\delta\) are moderate.

\begin{proposition}[Computer-assisted]\label{prop:witness}
Let \(2\le n\le7\), and let \((r,\delta)\) be \((2,\tfrac12)\), \((4,\tfrac12)\),
\((4,\tfrac12)\), \((8,\tfrac14)\), \((8,\tfrac14)\), \((8,\tfrac14)\) for \(n=2,\dots,7\). Let
\(x_j=3+\delta\,(j-n+1)\) for \(0\le j\le2n-2\), let the geometry be given by
\eqref{eq:normalize}, and let the kinetics be given by \eqref{eq:freedom} with
\(b_i=\beta_i=1\) and \(\theta_i=1/Y_i(X_\circ)\), where \(X_\circ\) is the coalesced
equilibrium \eqref{eq:coalstate} for the same \(r\); no relaxation \eqref{eq:relaxation} is
applied. Then the system with totals \((2r,2,2r+2)\) has exactly \(2n-1\) positive equilibria;
the \(n\) equilibria with even index are sinks, and at each of the \(n-1\) equilibria with odd
index exactly one eigenvalue of \(\JS\) has positive real part and none lies on the imaginary
axis.
\end{proposition}
\begin{proof}
For each of the \(48\) equilibria the rational matrix \(\JS\) was formed, its characteristic
polynomial was computed exactly by the Faddeev--LeVerrier recursion in integer arithmetic, with
the Cayley--Hamilton identity as a check, and the Routh array \cite[Ch.~XV]{Gantmacher} was
formed in rational arithmetic. In every case all \(3n+1\) entries of the first column are
nonzero, and the number of sign changes is \(0\) for even \(j\) and \(1\) for odd \(j\). The
count of equilibria is Theorem~\ref{thm:count}.
\end{proof}

\begin{example}[Two sites]\label{ex:two}
For \(n=2\), \(r=2\) and \(x=(\tfrac52,3,\tfrac72)\) one finds \(N=\tfrac{111}{4}+8u\),
\(J=\tfrac{15}{4}+32u\), \(D=1+\tfrac{8672}{2625}u\), \(B=\tfrac{1287}{125}+\tfrac{2288}{875}u\),
the totals \((\ET,\FT,\ST)=(4,2,6)\), and the rate constants
\[
\begin{array}{c|cccccc}
 i & a_i & b_i & c_i & \alpha_i & \beta_i & \gamma_i\\\hline
 1 & \tfrac{19782}{1375} & 1 & \tfrac{625}{1573} & \tfrac{147000}{124267} & 1 & \tfrac{45}{11}\\[3pt]
 2 & \tfrac{311808}{192049} & 1 & \tfrac{4065}{2431} & \tfrac{79}{34} & 1 & \tfrac{45}{34}
\end{array}
\]
The three equilibria, in the order \((S_0,S_1,S_2;\,E,F;\,C_1,C_2;\,Y_1,Y_2)\), are
\begin{align*}
 X_0&=\bigl(\tfrac{250}{693},\tfrac{1027}{1008},\tfrac{271}{528};\ \tfrac34,\tfrac87;\
 \tfrac{39}{14},\tfrac{13}{28};\ \tfrac{125}{462},\tfrac{271}{462}\bigr),\\
 X_1&=\bigl(\tfrac{2625}{11297},1,\tfrac{8672}{11297};\ 1,1;\
 \tfrac{189}{79},\tfrac{48}{79};\ \tfrac{2625}{11297},\tfrac{8672}{11297}\bigr),\\
 X_2&=\bigl(\tfrac{350}{2223},\tfrac{869}{912},\tfrac{4065}{3952};\ \tfrac54,\tfrac89;\
 \tfrac{77}{38},\tfrac{55}{76};\ \tfrac{875}{4446},\tfrac{1355}{1482}\bigr).
\end{align*}
\(X_0\) and \(X_2\) are sinks and \(X_1\) is a saddle. The eigenvalues of largest real part,
computed in floating point, are approximately \(-0.0120\), \(+0.0050\) and \(-0.0084\); the
remaining ones have real parts below \(-2\).
\end{example}

\subsection{A three-site benchmark with separated states}\label{sec:benchmark}
The witnesses above are close to coalescence. For a benchmark with well-separated states we
take \(n=3\),
\[
 x=(2,\tfrac72,5,\tfrac{13}{2},8),\qquad r=9,\qquad (\ET,\FT,\ST)=(18,2,20),
\]
for which \(N=\tfrac{10165}{4}+1926u+64u^2\), \(J=\tfrac{385}{2}+3562u+768u^2\) and
\[
 D=1+\tfrac{505005048}{38591245}u+\tfrac{48362752}{38591245}u^2,\qquad
 B=\tfrac{111414528}{1102607}+\tfrac{5946508756}{38591245}u+\tfrac{236379904}{38591245}u^2 .
\]
The nine kinetic parameters \((\theta_i,b_i,\beta_i)\) were chosen by a bounded numerical
search that maximizes the smallest decay rate at the three even-indexed equilibria, subject to
the ceilings \(a_i,\alpha_i\le2\) and \(b_i,c_i,\beta_i,\gamma_i\le2\) in model units; the
ceilings remove the trivial improvement obtained by rescaling time. The result was rounded to
rational numbers, and everything stated below refers to these exact rational rate constants,
which are distributed with the source. Three root geometries were compared in this way; the
smallest decay rates found were \(0.80\times10^{-4}\), \(0.98\times10^{-4}\) and
\(1.26\times10^{-4}\) for \(x=(2,3,4,5,6)\), \((\tfrac32,\tfrac52,\tfrac72,\tfrac92,\tfrac{11}2)\) and
the geometry selected here. These numbers are floating-point diagnostics, and the search makes
no claim of optimality.

\begin{proposition}[Computer-assisted]\label{prop:benchmark}
For the benchmark rate constants the class with totals \((18,2,20)\) contains exactly five
positive equilibria, listed in Table~\ref{tab:states}. The equilibria \(X_0,X_2,X_4\) are
sinks, and \(X_1,X_3\) are hyperbolic saddles with one unstable direction.
\end{proposition}
\begin{proof}
The rate constants realize the geometry above through \eqref{eq:freedom}, so
Theorem~\ref{thm:count} applies (\(r=9>\max_ju_j=\tfrac{63}{8}\) and
\(4r+\sqrt{1+8r}>36>25=\sum_jx_j\)). The stability statements are exact Routh computations as
in Proposition~\ref{prop:witness}.
\end{proof}

\begin{table}[ht]
\centering
\caption{The five equilibria of the three-site benchmark. Concentrations in the last two
columns are in \(\mu\)M for the unit \(C_0=0.1\,\mu\)M; decimal entries are rounded values of
exact rational numbers. ``Unstable'' is the number of eigenvalues of \(\JS\) in the open right
half-plane, from an exact Routh array.}\label{tab:states}
\begin{tabular}{cccccccc}
\toprule
\(j\) & \(x_j\) & \(u=E/F\) & \(E\) & \(F\) & \(\Phi'(u_j)\) & \(S_3+Y_3\) (\(\mu\)M) & unstable\\
\midrule
0 & \(2\) & \(3/8\) & \(1/2\) & \(4/3\) & \(-0.29235\) & \(0.003380\) & 0\\
1 & \(7/2\) & \(45/32\) & \(5/4\) & \(8/9\) & \(\phantom{-}0.02050\) & \(0.026743\) & 1\\
2 & \(5\) & \(3\) & \(2\) & \(2/3\) & \(-0.00772\) & \(0.072950\) & 0\\
3 & \(13/2\) & \(165/32\) & \(11/4\) & \(8/15\) & \(\phantom{-}0.01002\) & \(0.138019\) & 1\\
4 & \(8\) & \(63/8\) & \(7/2\) & \(4/9\) & \(-0.08439\) & \(0.216157\) & 0\\
\bottomrule
\end{tabular}
\end{table}

\begin{figure}[ht]
\centering
\includegraphics[width=\linewidth]{figures/benchmark.pdf}
\caption{The three-site benchmark. (a) The residual \(\Phi(u)-\ST\) on an
\(\operatorname{arcsinh}\) scale; filled markers are sinks, open markers are saddles. (b) The
readout \(S_3+Y_3\) at the five equilibria; the shaded bars are the certified readout intervals
of Proposition~\ref{prop:certificate}, including a measurement error of \(\pm0.001\,\mu\)M. The
curves and markers are drawn from exact rational data; the counts are certified by the
formulas and the Routh arrays, not by the plot.}\label{fig:benchmark}
\end{figure}

\paragraph{Declared units.}
To read the benchmark in physical units we declare a concentration unit \(C_0=0.1\,\mu\)M and a
time unit \(T_0=10\) s. Bimolecular constants are divided by \(C_0T_0\) and unimolecular
constants by \(T_0\). The totals become \((\ET,\FT,\ST)=(1.8,\,0.2,\,2.0)\,\mu\)M, and the rate
constants are those of Table~\ref{tab:rates}. The association constants are at most
\(2\,\mu\mathrm M^{-1}\mathrm s^{-1}\) and the unimolecular constants at most
\(0.2\,\mathrm s^{-1}\). The slowest of the three numerically computed decay rates corresponds
to an e-folding time of about \(22\) hours. These units are declarations, not fits: the
unimolecular constants span a factor of about \(33{,}446\), the kinase total is comparable to
the substrate total, and no measured kinetic data were used. The benchmark is a reproducible
test case, not a model of a particular kinase--phosphatase pair.

\begin{table}[ht]
\centering
\caption{Rate constants of the three-site benchmark in the declared units, rounded to six
significant digits. The exact rational values in model units are distributed with the
source and are authoritative.}\label{tab:rates}
\begin{tabular}{ccccccc}
\toprule
Site & \(a_i\) (\(\mu\mathrm M^{-1}\mathrm s^{-1}\)) & \(b_i\) (\(\mathrm s^{-1}\)) &
\(c_i\) (\(\mathrm s^{-1}\)) & \(\alpha_i\) (\(\mu\mathrm M^{-1}\mathrm s^{-1}\)) &
\(\beta_i\) (\(\mathrm s^{-1}\)) & \(\gamma_i\) (\(\mathrm s^{-1}\))\\
\midrule
1 & \(2.00000\) & \(5.96184\times10^{-6}\) & \(0.00197333\) & \(0.171890\) & \(0.0427272\) & \(0.199398\)\\
2 & \(1.86875\) & \(7.83183\times10^{-4}\) & \(0.0162999\) & \(1.99626\) & \(0.0268102\) & \(0.191934\)\\
3 & \(0.217843\) & \(0.0126769\) & \(0.0383203\) & \(2.00000\) & \(0.0127038\) & \(0.187296\)\\
\bottomrule
\end{tabular}
\end{table}

\subsection{A rigorous local certificate for independent parameter errors}
Theorem~\ref{thm:main}(c) gives an open set of parameters without an estimate of its size. For
the benchmark we certify an explicit, if very small, neighbourhood, in which all \(18\) rate
constants and all three totals vary independently; in particular the perturbed systems need
not lie in the source-preserving family \eqref{eq:freedom}, and their equilibria move.

Let \(q=(S_1,S_2,S_3,C_1,C_2,C_3,Y_1,Y_2,Y_3)\) be the coordinates on the class, let \(q_0\) be
one of the three nominal sinks, and use the weighted deviations
\(e=\diag(q_0)^{-1}(q-q_\theta)\) from the equilibrium \(q_\theta\) of a perturbed parameter
vector \(\theta\). Let \(\JS_0\) be the restricted Jacobian at \(q_0\) in the weighted
coordinates.

\begin{proposition}[Computer-assisted]\label{prop:certificate}
Let \(\epsilon_{\mathrm{rel}}=1/671088640000\approx1.49\times10^{-12}\), and let each of the
\(18\) rate constants and the three totals of the benchmark vary independently within the
relative distance \(\epsilon_{\mathrm{rel}}\) of its nominal value. For each of the three
nominal sinks \(X_0,X_2,X_4\), with the radii \(\varrho_0\) and \(\varrho_{\mathrm{in}}\) of
Table~\ref{tab:certificate}:
\begin{enumerate}
\item[\textup{(i)}] the perturbed system has exactly one equilibrium \(q_\theta\) in the box
\(\|\diag(q_0)^{-1}(q-q_0)\|_\infty\le\varrho_0\); it is positive and satisfies
\(\|\diag(q_0)^{-1}(q_\theta-q_0)\|_\infty\le\varrho_0/4\);
\item[\textup{(ii)}] there is a rational symmetric positive definite matrix \(\bP\), the same
for all admitted parameters, such that every solution with \(\|e(0)\|_2\le\varrho_{\mathrm{in}}\)
stays in that box and satisfies
\[
 \|e(t)\|_2\le\frac{\varrho_0}{2}\,\exp\Bigl(-\frac{c}{2M}\,\frac{t}{T_0}\Bigr),
 \qquad M=\|\bP\|_\infty,\quad c=\frac{1}{\|\bQ^{-1}\|_\infty}-\ell_{\mathrm{loss}}>0,
\]
where \(\bQ=-(\JS_0^{T}\bP+\bP\JS_0)\) and \(\ell_{\mathrm{loss}}\) is an interval bound for the
variation of \(\JS^{T}\bP+\bP\JS\) over the box and over all admitted parameters; in particular
\(\|e(t)\|_2\) has decreased below \(\varrho_0/20\) after the time \(T_{10}=60M/c\) seconds of
Table~\ref{tab:certificate};
\item[\textup{(iii)}] the readouts \(S_3+Y_3\) at the three perturbed sinks, with a measurement
error of \(\pm0.001\,\mu\mathrm M\) added, lie in the pairwise disjoint intervals
\([0.00237970,0.00437971]\), \([0.07195005,0.07395006]\) and
\([0.21515666,0.21715672]\) (\(\mu\)M).
\end{enumerate}
\end{proposition}
\begin{proof}
All steps are carried out in rational or rational-interval arithmetic. A floating-point
Lyapunov solve proposes \(\bP\), which is then rounded to a rational matrix; exact \(LDL^{T}\)
factorizations verify that \(\bP\) and \(\bQ\) are positive definite. Since \(\bP\) and
\(\bQ^{-1}\) are symmetric, their spectral norms are at most their \(\infty\)-norms, so
\(\bP\preceq M\bI\) and \(\bQ\succeq\|\bQ^{-1}\|_\infty^{-1}\bI\).

(i) For the Newton-type map \(q\mapsto q-\JS_0^{-1}g_\theta(q)\) (in weighted coordinates),
interval evaluation over the box and the parameter set shows that its derivative has
\(\infty\)-norm at most some \(\vartheta\le\tfrac12\) and that it moves \(q_0\) by at most
\((1-\vartheta)\varrho_0/4\); by the contraction principle it has a unique fixed point in the
box, within \(\varrho_0/4\) of \(q_0\), and the fixed points of this map in the box are exactly
the equilibria there. The same interval evaluation bounds all twelve concentrations away from zero. Here it
matters that each reaction vector is multiplied by \(\JS_0^{-1}\) (and, below, by \(\bP\))
\emph{before} the interval evaluation: this keeps the exact cancellations between the
reactions and improves the admissible \(\epsilon_{\mathrm{rel}}\) by about eight orders of
magnitude over a product of separate norm bounds.

(ii) For admitted \(\theta\) and \(q\) in the box, the segment from \(q_\theta\) to \(q\) lies in
the box, and the mean-value form of the vector field gives
\(\dot e=\bar\JS e\), where \(\bar\JS\) is an average of restricted Jacobians over the segment. The
interval bound gives
\(\|(\bar\JS-\JS_0)^{T}\bP+\bP(\bar\JS-\JS_0)\|_2\le\ell_{\mathrm{loss}}\), so for
\(V=e^{T}\bP e\),
\[
 \dot V\le-\bigl(\|\bQ^{-1}\|_\infty^{-1}-\ell_{\mathrm{loss}}\bigr)\|e\|_2^2=-c\,\|e\|_2^2
 \le-\frac cM\,V .
\]
Let \(p_{\min}\) be a rational lower bound for the smallest eigenvalue of \(\bP\). The
ellipsoid \(V\le p_{\min}(\varrho_0/2)^2\) lies in the Euclidean ball of radius \(\varrho_0/2\)
around \(q_\theta\), hence in the box, because the centre has moved by at most
\(\varrho_0/4\) in the \(\infty\)-norm; it is invariant; and it contains the ball
\(\|e\|_2\le\varrho_{\mathrm{in}}:=(\varrho_0/2)\sqrt{p_{\min}/M}\) (rounded down). On it,
\(\|e(t)\|_2^2\le V(t)/p_{\min}\le(\varrho_0/2)^2\exp(-ct/M)\) in model time. Since
\(\log10<3\), a tenfold decrease of this envelope is reached by the model time \(6M/c\), that
is, by \(60M/c\) seconds.

(iii) The readout is a linear function of \(q\), and (i) bounds the displacement of each sink.
\end{proof}

\begin{table}[ht]
\centering
\caption{The local certificate of Proposition~\ref{prop:certificate}. The column ``earlier''
is the bound obtained when the norms are rounded to powers of ten and half of the Lyapunov
margin is discarded; the column \(T_{10}\) keeps the exact constants. Times are rounded
upwards.}\label{tab:certificate}
\begin{tabular}{ccccccc}
\toprule
Sink & \(\varrho_0\) & \(\varrho_{\mathrm{in}}\) & \(M\) & \(c\) & \(T_{10}\) (s) & earlier (s)\\
\midrule
\(X_0\) & \(1/8192000\) & \(1/16384000000\) & \(5846.11\) & \(0.98516\) & \(357{,}000\) & \(1.2\times10^{7}\)\\
\(X_2\) & \(1/16384000\) & \(1/327680000000\) & \(535109.89\) & \(0.97537\) & \(32{,}918{,}000\) & \(1.2\times10^{9}\)\\
\(X_4\) & \(1/2048000\) & \(1/40960000000\) & \(383168.49\) & \(0.98112\) & \(23{,}433{,}000\) & \(1.2\times10^{9}\)\\
\bottomrule
\end{tabular}
\end{table}

The certificate should be read for what it is. The parameter tolerance
\(\epsilon_{\mathrm{rel}}\approx1.5\times10^{-12}\) and the neighbourhoods
\(\varrho_{\mathrm{in}}\) are far too small to represent laboratory tolerances, and the times
\(T_{10}\) are upper bounds for an envelope inside an invariant ellipsoid, not predicted
recovery times: the worst bound is about \(381\) days, whereas the numerically computed slowest
e-folding time is about \(22\) hours. Much of the gap between the two reflects the conditioning of
\(\bP\) (\(M\approx5\times10^5\)) rather than the dynamics. Keeping the exact rational
constants instead of rounded ones improved the worst bound \(36\)-fold, from about \(38\)
years, without changing the system, the parameter box or the neighbourhoods. What the
proposition does establish is that a full-dimensional, explicitly described set of parameters
supports three distinguishable, locally attracting states with quantified return.

\section{Formal verification and reproducibility}\label{sec:formal}
The development in Lean~4 \cite{Lean4} with Mathlib \cite{Mathlib} has two parts. All modules
compile without warnings under Lean \texttt{v4.30.0} with warnings treated as errors, contain no
\texttt{sorry}, and depend only on the axioms \texttt{propext}, \texttt{Classical.choice} and
\texttt{Quot.sound}.

\emph{The equilibrium count.} The modules \texttt{Recurrence}, \texttt{Conversion},
\texttt{Source}, \texttt{Assembly} and \texttt{Resolution} of the namespace
\texttt{PhosphorylationSharpness} define the state space, the six rate constants per site, the
mass-action vector field of \eqref{eq:network} componentwise, and the three totals, and prove
\begin{quote}
\texttt{theorem sharp\_lower\_bound (n :\ $\mathbb N$) (hn :\ 1 $\le$ n) :\ Attains n (2*n-1)}
\end{quote}
where \texttt{Attains n k} states that there exist positive rate constants, positive totals and
an injective family of \(k\) positive states at which every component of the vector field
vanishes and the three totals take the given values. The statement contains no reduced
equation and no numerical input. The proof follows Section~\ref{sec:construction} for the roots
\(x_j=j+2\), except that positivity of \(D_{\raw}\) and \(B_{\raw}\) is proved for all
sufficiently large \(r\) by an elementary estimate based on \eqref{eq:qformula} instead of
Theorem~\ref{thm:positivity}. A further module, \texttt{PublicationAlgebra}, verifies
\eqref{eq:product} and coefficient positivity for arbitrary lists of roots greater than one,
the determinant \eqref{eq:pivot}, the factorization \eqref{eq:residual}, and the sign of the
prefactor in \eqref{eq:slope}.

\emph{The algebra of the stability proof.} The namespace \texttt{PhosphorylationStableCapacity}
contains four modules. \texttt{StructuralAlgebra} proves the coalesced identity
\((x-1)N-2J=(x-3)^{2m+1}\) for every \(m\) (\texttt{coalesced\_identity}); that the relaxation
\eqref{eq:relaxation} preserves the complex balances and positivity
(\texttt{relaxation\_balance}, \texttt{relaxation\_positive}); that a matrix with vanishing
last column has characteristic polynomial \(\lambda\) times that of its leading block
(\texttt{slow\_link\_charpoly}), as used in Lemma~\ref{lem:slowlink}; the elimination
\eqref{eq:alpha_beta}\(\Rightarrow\)\eqref{eq:feedbacklimit}
(\texttt{loaded\_feedback\_elimination}); the monotonicity step behind \(\phi_i>0\)
(\texttt{feedback\_ratio\_step}); the reduction of \(\bB_*\bz\) to the scalar gain
(\texttt{feedback\_vector}); and the identity and positivity \eqref{eq:gainfinal} for every
natural number \(n\ge2\) and every \(0<v<1\) (\texttt{feedback\_gain\_identity},
\texttt{feedback\_gap\_positive}, \texttt{feedback\_gain\_lt\_one}). \texttt{KineticFreedom}
proves, for the vector field of the first part, that the retuning \eqref{eq:freedom} keeps all
rate constants positive and preserves every equilibrium satisfying the three balance
relations (\texttt{retune\_positive}, \texttt{retune\_equilibrium}). \texttt{Assembly} combines
these into the statement that, given the matrix identities \(\bH\bz=\one\) and the value of
\(\bphi^{T}\bz\), every entry of \(\bB_*\bz\) is negative
(\texttt{large\_r\_block\_certificate}). \texttt{DiagonalAlgebra} verifies the pairwise identity
behind \eqref{eq:energy}, the congruence used for \(\bA_r\boldsymbol\Omega\), and the sign of
\(\bQ_*\bz\).

\emph{What is not formalized.} Theorem~\ref{thm:main} as a whole is not a formal theorem. The
following parts are proved only in the conventional way: the ceiling of
Section~\ref{sec:ceiling} (degree theory, transversality, and the quoted
Theorem~\ref{thm:WS}); the interlacing lemma, Theorem~\ref{thm:positivity} and
Lemma~\ref{lem:ratios}; the determinant and slope formulas; the derivation of \(\bM\), \(\bh\)
and \(\bK\) and the limits of Section~\ref{sec:prefix}, whose \emph{outputs} are the hypotheses
of the formal algebraic statements; the \(M\)-matrix and Gershgorin arguments; and all spectral
perturbation arguments of Section~\ref{sec:lift}. A compiled conditional identity does not
prove its analytic premises. Propositions~\ref{prop:witness}, \ref{prop:benchmark}
and~\ref{prop:certificate} are exact computations that are not formalized.

{\sloppy\emph{Reproduction.} The source package contains the Lean modules, the exact data of
the benchmark, and a Python script, \texttt{check\_paper.py}, which uses only the standard
library and rational arithmetic. It verifies the construction for \(n\le6\) (equilibrium
equations, totals, the identity \eqref{eq:freeinventory}) for distinct and for coalesced roots;
Lemma~\ref{lem:interlace} by Sturm sequences, Lemma~\ref{lem:ratios}, and
Theorem~\ref{thm:positivity} at the boundary of its range, on random root configurations with
and without repeated roots; the moments \eqref{eq:moments} and \eqref{eq:vmu}, the positivity
of \(\bphi\), the gain formula \eqref{eq:gainfinal}, Table~\ref{tab:deficit}, and the
definiteness of \(\bQ_*\), for \(n\le8\); the exact identities of Lemma~\ref{lem:compression},
the singularity of \(\bK\), and the \(O(1/r)\) convergence \(r\bA_r\to\bB_*\) for \(n\le5\); the
determinant formula \eqref{eq:det} and the readout formula \eqref{eq:readout}; the base case
\(n=1\); Proposition~\ref{prop:witness} (for \(n\le4\) by default and for \(n\le7\) with the
option \texttt{-{}-full}); Example~\ref{ex:two}; and every number in
Section~\ref{sec:benchmark}, including an independent replay of the interval bounds of
Proposition~\ref{prop:certificate}. These computations test the exposition and the
implementation. The proofs of Theorems~\ref{thm:main} and~\ref{thm:split} do not depend on
them.\par}

\section{Discussion}\label{sec:discussion}
The sequential distributive \(n\)-site cycle has equilibrium capacity \(2n-1\) and stable-state
capacity \(n\), for every \(n\). Both values are attained by one explicit rational system, and
the stable ceiling holds without any nondegeneracy assumption. The reusable idea in the proof
is the separation of geometry and kinetics: the equilibrium equations see only half of the
parameters, so one can first build the richest possible equilibrium set, then coalesce it, and
then use the remaining parameters to organize time scales so that the transverse spectrum
becomes computable. In the present case the computation ends in a scalar gain with a closed
form.

We close with limitations and open questions.

\emph{Conditioning.} The proof is an existence proof with ordered limits. The margin
\eqref{eq:deficit} decays like \(n^2(2/9)^n\), the splitting law \eqref{eq:splitlaw} shows that
near coalescence the slowest rate is of order \(\delta^{2n-2}\), and the open set of
Theorem~\ref{thm:main}(c) comes without a size. Proposition~\ref{prop:witness} and the
benchmark show that for small \(n\) much better constants are available than the proof
suggests. It would be interesting to find a construction of \(n\) sinks whose spectral gap and
parameter tolerance decay at most polynomially in \(n\), or to show that none exists.

\emph{Enzyme levels.} In our examples \(\ST=\ET+\FT\) and the kinase is in large excess over
the phosphatase, so that most of the substrate is bound to the kinase; this is far from the
substrate-excess regime of Michaelis--Menten reductions, in which generically at most \(n+1\)
equilibria exist \cite[Theorem~4]{WangSontag2008}. How small \(\ET/\ST\) and \(\FT/\ST\) can be
in a system with \(2n-1\) equilibria, or with \(n\) sinks, is open. By the symmetry
\(E\leftrightarrow F\), \(S_i\leftrightarrow S_{n-i}\) of \eqref{eq:network} all results hold
with the roles of the enzymes exchanged.

\emph{Global dynamics.} Theorem~\ref{thm:main} bounds the number of attracting equilibria, not
the number of attractors. For \(n=1\) every trajectory converges to the unique equilibrium
\cite{AngeliSontag2008}. Whether the two-site cycle admits periodic orbits is a
well-known open question; see \cite{CMS2019,CFM2020,ConradiMincheva2024} for partial results
and related mechanisms. Our ceiling neither uses nor implies the absence of oscillations for
any \(n\), and attracting periodic orbits, if present, are not counted by \(\capC\). We do not
know the basins of the \(n\) sinks constructed here, and we do
not claim that they attract almost all initial conditions.

\emph{Thermodynamics and noise.} The catalytic steps in \eqref{eq:network} are irreversible
effective reactions. The constructed states carry a positive phosphorylation flux and
therefore presuppose a maintained chemical driving force; a closed system with a finite supply
of ATP cannot hold them indefinitely. A model with explicit fuel, or with reversible catalytic
steps, is a different network and requires its own analysis. Likewise, the theorem concerns
the deterministic rate equations. Whether the \(n\) states are retained for a useful time at
finite copy numbers depends on the volume, the readout and the time horizon, and neither
retention nor its failure is established here. For these reasons Corollary~\ref{cor:labels}
should be read as a statement about stable equilibrium labels of a mechanism, not as a claim
that a cell realizes or inherits them.

\emph{Other mechanisms.} The method uses the monomial parametrization of the equilibria, a
locus on which the scalar equation factors after a quadratic change of variable, and the path
structure of the slow dynamics. It is natural to ask for analogues in cycles with several
kinases or phosphatases, in mixed processive--distributive mechanisms, and in cascades, for
which sharp bounds are not known. The ceiling of Section~\ref{sec:ceiling} applies verbatim to
any mass-action network with compact compatibility classes and without boundary equilibria
for which the number of positive equilibria in a class is at most \(2k-1\) for all rate
constants: such a network has at most \(k\) locally asymptotically stable equilibria in each
class.
